\chapter{Especificación de Requisitos Software}
\label{cap:SRS}

Versión 0.1\\

\section{Introducción}

\subsection{Propósito del documento}

El propósito del presente documento es servir como especificación de requisitos de software (SRS) para las funcionalidades de estadísticas y logros del juez \textit{online} \aer. Este documento constituye la base para el diseño, desarrollo, validación y mantenimiento de dichas funcionalidades.

El documento está dirigido principalmente a los desarrolladores y mantenedores de las funcionalidades descritas, a los mantenedores del juez \aer, y a cualquier otra persona interesada en comprender el funcionamiento, alcance y limitaciones del producto.

Asimismo, este documento debe emplearse como referencia durante el ciclo de vida del software para constatar las funcionalidades implementadas, las decisiones de diseño adoptadas, los motivos que las justifican y las limitaciones actuales del sistema, con el objetivo de facilitar su mantenibilidad y evolución futura.

\subsection{Ámbito del producto}

El producto no cuenta con una denominación propia. A lo largo de este documento se hará referencia a él como \textit{las funcionalidades de estadísticas y logros para el juez ¡Acepta el reto!}, manteniendo la consistencia con el título del proyecto en cuyo marco se desarrolla.

En este documento se describen las especificaciones correspondientes a la primera versión del producto. Las capacidades principales incluyen la visualización fiel y actualizada de estadísticas de usuario, estadísticas asociadas a los ejercicios, y la gestión de logros o medallas otorgadas por la realización de determinadas acciones.

Cabe destacar que este documento se refiere exclusivamente a las funcionalidades de estadísticas y logros, las cuales se implementan de manera independiente a la arquitectura actual de \aer. No obstante, el producto se diseña con el objetivo de facilitar una posible integración futura en la arquitectura del juez, en caso de que se considere oportuno.

\subsection{Definiciones, acrónimos y abreviaciones}

\begin{itemize}
    \item \textbf{UI}: User Interface.
    \item \textbf{SRS}: Software Requirements Specification.
    \item \textbf{TFG}: Trabajo de Fin de Grado.
    \item \textbf{¡Acepta el reto!}: Juez \textit{online} de programación competitiva objeto de este proyecto.
\end{itemize}

\subsection{Referencias}

\begin{itemize}
    \item IEEE Std 830-1998 / IEEE Std 29148-2018, \textit{Software Requirements Specification}.
\end{itemize}

\section{Descripción general del producto}

\subsection{Perspectiva del producto}

El producto se desarrolla en el marco de un TFG, bajo la supervisión de los creadores y mantenedores de \aer. Su objetivo es cubrir la necesidad de introducir mayor dinamismo y nuevas fuentes de motivación para los usuarios del juez \textit{online} mediante estadísticas detalladas y un sistema de logros.

Aunque la integración directa de estas funcionalidades en el juez no forma parte del alcance del proyecto, el diseño tendrá en cuenta criterios de modularidad y desacoplamiento que faciliten una posible integración futura.

\subsection{Funciones del producto}

\begin{itemize}
    \item Mostrar estadísticas agregadas y detalladas de los usuarios que reflejen su progreso y actividad.
    \item Mostrar estadísticas de ejercicios que permitan inferir su dificultad, popularidad o tasa de éxito.
    \item Actualizar las estadísticas de manera automática y en tiempo cercano al real.
    \item Otorgar logros o medallas a los usuarios como recompensa por acciones específicas o hitos alcanzados.
    \item Ordenar a los usuarios en una tabla de clasificación en base a su puntuación.
\end{itemize}

\subsection{Limitaciones del producto}

El producto estará sujeto a limitaciones técnicas derivadas del \textit{stack} tecnológico seleccionado durante el desarrollo del TFG, así como a restricciones de tiempo y recursos propios de un proyecto académico. Asimismo, el sistema dependerá de la disponibilidad y calidad de los datos proporcionados por el juez \aer u otras fuentes externas. No se contemplan requisitos específicos de hardware más allá de los necesarios para ejecutar una aplicación web moderna.

\subsection{Características de los usuarios}

Se identifican los siguientes grupos de usuarios:

\begin{itemize}
    \item \textbf{Usuarios finales}: estudiantes o programadores que utilizan el juez. Tienen conocimientos técnicos básicos, acceden con frecuencia al sistema y su objetivo principal es medir su progreso y obtener motivación adicional para mejorar sus habilidades de programación y algoritmia.
    \item \textbf{Administradores/mantenedores}: usuarios con conocimientos técnicos avanzados y acceso a funciones de configuración o mantenimiento. Su uso es esporádico y orientado a la supervisión del sistema.
\end{itemize}

No se prevén necesidades especiales de accesibilidad en esta primera versión, aunque se seguirán buenas prácticas de diseño de interfaces.

\subsection{Supuestos y dependencias}

Se asume la disponibilidad de datos fiables sobre usuarios y ejercicios, así como el acceso a interfaces o mecanismos que permitan obtenerlos. El producto depende de librerías y servicios de terceros para la persistencia de datos y la visualización de estadísticas. Si alguna de estas dependencias no se cumple, el impacto principal será la imposibilidad de mostrar estadísticas actualizadas o de calcular correctamente los logros.

\section{Requisitos} 
\label{sec:requisitos}

\subsection{Interfaces externas}

El sistema expone y consume las siguientes interfaces externas:
\begin{itemize}
    \item \textbf{Interfaces de usuario (UI)}
 El sistema proporcionará una interfaz web accesible mediante navegadores modernos (Chrome, Firefox, Edge, etc).
 La UI incluirá:
    \begin{itemize}
        \item Paneles de estadísticas de usuario.
        \item Vistas de estadísticas de ejercicios.
        \item Sección de logros y medallas obtenidas.
        \item Tabla de clasifiación de usuarios.

 Se seguirán principios de diseño responsivo y consistencia visual con el estilo general de \aer, aunque modernizando o alterando ciertos elementos.
    \end{itemize}

    \item \textbf{Interfaces de hardware}
 No se definen interfaces hardware específicas. El sistema se ejecutará sobre infraestructura estándar de servidores web y será accesible desde dispositivos de escritorio y móviles.
\end{itemize}

\subsection{Funcionales}
\label{subsec:funcionales}

\textbf{REQ-FUN-001}
\begin{itemize}
    \item \textbf{Título:} Visualización de estadísticas de usuario
    \item \textbf{Descripción:} El sistema debe permitir visualizar estadísticas asociadas a un usuario, incluyendo número de ejercicios resueltos, envíos realizados y tasa de aciertos, entre otros.
    \item \textbf{Racional:} Proporcionar al usuario información clara sobre su progreso y actividad.
    \item \textbf{Criterio de aceptación:} Dado un usuario válido, el sistema muestra sus estadísticas de forma correcta y consistente.
    \item \textbf{Método de verificación:} Test funcional.

\end{itemize}

\textbf{REQ-FUN-002}
\begin{itemize}
    \item \textbf{Título:} Visualización de estadísticas de ejercicios
    \item \textbf{Descripción:} El sistema debe mostrar estadísticas agregadas por ejercicio, como número de intentos, número de resoluciones y porcentaje de éxito.
    \item \textbf{Racional:} Ayudar a los usuarios a estimar la dificultad o popularidad de los ejercicios.
    \item \textbf{Criterio de aceptación:} Las estadísticas mostradas coinciden con los datos almacenados.
    \item \textbf{Método de verificación:} Test funcional.

\end{itemize}

\textbf{REQ-FUN-003}
\begin{itemize}
    \item \textbf{Título:} Actualización automática de estadísticas y puntos
    \item \textbf{Descripción:} El sistema debe actualizar las estadísticas y puntos de usuario cuando se produzcan nuevos eventos relevantes (por ejemplo, un envío evaluado).
    \item \textbf{Racional:} Garantizar que la información mostrada sea actual y fiable.
    \item \textbf{Criterio de aceptación:} Las estadísticas reflejan los cambios sin necesidad de intervención manual.
    \item \textbf{Método de verificación:} Demostración, test funcional.

\end{itemize}

\textbf{REQ-FUN-004}
\begin{itemize}
    \item \textbf{Título:} Gestión de logros
    \item \textbf{Descripción:} El sistema debe otorgar logros a los usuarios cuando se cumplan determinadas condiciones predefinidas.
    \item \textbf{Racional:} Incrementar la motivación de los usuarios mediante recompensas simbólicas.
    \item \textbf{Criterio de aceptación:} Un logro se asigna automáticamente cuando se cumple su condición.
    \item \textbf{Método de verificación:} Test funcional.


\end{itemize}

\textbf{REQ-FUN-005}
\begin{itemize}
    \item \textbf{Título:} Consulta de logros obtenidos
    \item \textbf{Descripción:} El sistema debe permitir a un usuario consultar la lista de logros que ha obtenido.
    \item \textbf{Racional:} Hacer visibles las recompensas y fomentar la participación.
    \item \textbf{Criterio de aceptación:} El usuario puede ver todos sus logros sin errores.
    \item \textbf{Método de verificación:} Demostración y test funcional.
\end{itemize}

\textbf{REQ-FUN-006}
\begin{itemize}
    \item \textbf{Título:} Actualización de la tabla de clasificación
    \item \textbf{Descripción:} El sistema debe actualizar automáticamente la tabla de clasificación según las puntuaciones de los usuarios.
    \item \textbf{Racional:} Mostrar una imagen fiel y actualizada de las puntuaciones de los usuarios.
    \item \textbf{Criterio de aceptación:} La tabla de clasificación se actualiza correctamente cada vez que un usuario completa una acción que afecta su puntuación.
    \item \textbf{Método de verificación:} Test funcional.
\end{itemize}

\textbf{REQ-FUN-007}
\begin{itemize}
    \item \textbf{Título:} Consulta de la tabla de clasificación
    \item \textbf{Descripción:} El sistema debe permitir a los usuarios consultar la tabla de clasificación, mostrando su posición y la de otros usuarios según su puntuación.
    \item \textbf{Racional:} Permitir a los usuarios compararse con otros.
    \item \textbf{Criterio de aceptación:} El usuario puede visualizar la tabla de clasificación completa, incluyendo su posición, sin errores y con los datos actualizados.
    \item \textbf{Método de verificación:} Demostración y test funcional.
\end{itemize}

\textbf{REQ-FUN-008}
\begin{itemize}
    \item \textbf{Título:} Compartición de estadísticas y logros
    \item \textbf{Descripción:} El sistema debe permitir a los usuarios compartir sus estadísticas o logros de usuario como imágen.
    \item \textbf{Racional:} Permitir a los usuarios compartir sus hitos, generando más motivación.
    \item \textbf{Criterio de aceptación:} El usuario puede recibir desde la aplicación una imágen o link con la que mostrar una estadística o logro recibido a otras personas.
    \item \textbf{Método de verificación:} Demostración.
\end{itemize}

\subsection{No funcionales}

\textbf{REQ-NF-001}
\label{req:REQ-NF-001}
\begin{itemize}
    \item \textbf{Título:} Rendimiento
    \item \textbf{Descripción:} El sistema debe responder a las consultas de estadísticas en menos de 2 segundos en condiciones normales de carga.
    \item \textbf{Racional:} Garantizar una experiencia de usuario fluida.
    \item \textbf{Criterio de aceptación:} El tiempo de respuesta se mantiene dentro del límite establecido.
    \item \textbf{Método de verificación:} Test de rendimiento.
\end{itemize}

\textbf{REQ-NF-002}
\label{req:REQ-NF-002}
\begin{itemize}
    \item \textbf{Título:} Fiabilidad de los datos
    \item \textbf{Descripción:} Las estadísticas mostradas deben ser consistentes con los datos de origen.
    \item \textbf{Racional:} Evitar información incorrecta que afecte a la confianza del usuario.
    \item \textbf{Criterio de aceptación:} No se detectan discrepancias en auditorías de datos.
    \item \textbf{Método de verificación:} Análisis? TEST FUNCIONAL?.
\end{itemize}

\textbf{REQ-NF-003}
\label{req:REQ-NF-003}
\begin{itemize}
    \item \textbf{Título:} Usabilidad
    \item \textbf{Descripción:} La interfaz debe ser comprensible sin necesidad de formación previa.
    \item \textbf{Racional:} Facilitar la adopción del sistema por parte de usuarios sin experiencia.
    \item \textbf{Criterio de aceptación:} Usuarios de prueba completan tareas básicas sin ayuda.
    \item \textbf{Método de verificación:} Test de usabilidad.
\end{itemize}

\textbf{REQ-NF-004}
\label{req:REQ-NF-004}
\begin{itemize}
    \item \textbf{Título:} Capacidad de integrarse con el OJ \aer
    \item \textbf{Descripción:} El diseño e implementación de la aplicación deben posibilitar la futura integración de la aplicación en la estructura general de \aer, gracias a un enfoque modular, y comprensión del sistema de eventos existente.
    \item \textbf{Racional:} Facilitar la puesta en marcha del sistema en el entorno real (el OJ \aer), donde podrá ofrecer su utilidad completa.
    \item \textbf{Criterio de aceptación:} La aplicación debe demostrar una arquitectura modular, con una capa de negocio independiente de la interfaz y el manejo de la base de datos, y usar sistemas de manejo de eventos abstractos que se pueda reutilizar y cambiar fácilmente. 
    \item \textbf{Método de verificación:} Revisión del código en grupo para validar el cumplimiento de los principios de diseño modular y desacoplado.
\end{itemize}

\subsection{Diseño e implementación}

El sistema deberá diseñarse siguiendo principios de modularidad y desacoplamiento para facilitar su mantenimiento y posible integración futura.

 Se priorizará el uso de tecnologías web estándar y ampliamente soportadas.

 El sistema deberá ser portable entre navegadores modernos y fácilmente desplegable en entornos de servidor convencionales.

 Las decisiones de diseño deberán documentarse para facilitar la mantenibilidad y reutilización del código.

\section{Verificación}

La verificación del sistema se llevará a cabo comprobando que cada requisito definido en la Sección 3 ha sido satisfecho de forma objetiva. Para ello, se emplearán distintos métodos de verificación, incluyendo pruebas funcionales, análisis y demostraciones.

Cada requisito estará asociado a uno o más casos de prueba que permitan validar su cumplimiento. Los resultados de la verificación se documentarán, indicando si el requisito ha sido cumplido, parcialmente cumplido o no cumplido, así como las incidencias detectadas y las acciones correctivas adoptadas.